\section{Sample and data}\label{sec:data}
% Objective: The calendar constraint from Rule 424(b)(2); the screening census; the six notes and the structures they span.

\subsection{Why the sample is one week}

A note can only be repriced against its disclosed value while its terms are
both fixed and public, and that state begins later than one might expect and is
announced only briefly. The preliminary supplement carries ranges: a coupon quoted as
``at least'' some rate, a trigger to be determined, an estimated value given as
an interval.\footnote{The phrasing is standard across the preliminary
supplements in this week. One example, filed under accession number
0001918704-26-031175, offers ``a Contingent Interest Rate of at least 9.70\%
per annum.'' These documents are superseded by their priced counterparts and
are not preserved in the repository.} The priced supplement replaces all of it with numbers, and Rule
424(b)(2) requires that document to be filed no later than the second business
day following the earlier of the date the offering price was determined or the
date the supplement is first used \citep{SEC424b2}.

So the terms arrive in a stream of final supplements with a two-day lag and no
subsequent amendment. Choosing a calendar week and taking every qualifying
filing in it is therefore a sampling rule a reader can check, and one that
cannot be tuned after the fact by extending the window until the results
improve. The week used here is 5 to 9 October 2026. Five of the six notes were
priced on 7 October and the sixth on 5 October.

\subsection{What the week contained}

% Generated by scripts/make_paper_tables.py from the frozen results.
% Do not edit by hand; the next build overwrites it.
\begin{table}[tbp]
\centering
\caption{Every 424B2 filed in the pricing week, by the reason it does or does not reach the sample}
\label{tab:census}
\small
\begin{tabular}{l r r}
\toprule
Disposition & Filings & Share (\%) \\
\midrule
Preliminary supplement; terms not yet fixed & 667 & 61.4 \\
Not a worst-of structure & 201 & 18.5 \\
No contingent coupon & 84 & 7.7 \\
No estimated value disclosed & 5 & 0.5 \\
Single-name or fund underlyings & 16 & 1.5 \\
An underlying with no listed options & 42 & 3.9 \\
Priceable from public option quotes & 71 & 6.5 \\
\midrule \textit{Total 424B2 filings} & \textit{1,086} & \textit{100.0} \\
\bottomrule
\end{tabular}
\notetext{\textit{Notes:} All 424B2 filings retrieved from EDGAR full-text search for the week of the snapshot. Each filing is assigned to the first stage it fails, so the rows partition the population. A preliminary supplement is excluded because its coupon, trigger and estimated value are still ranges. The last three rows are the final worst-of notes with a contingent coupon and a disclosed estimated value, 129 in all, of which 71 (55\%) can be priced from public option quotes; the six notes studied here are drawn from those 71. Source: \texttt{src/data/edgar\_notes.py}.}
\end{table}


Table~\ref{tab:census} is the whole filing population, each filing assigned to
the first stage it fails, so the rows partition rather than overlap. Of 1{,}086
supplements filed on Form 424B2 that week, 667 are preliminary and 285 are
final but not the structure under study, either because they are not worst-of or
because they carry no contingent coupon. A further five are final worst-of notes
with a contingent coupon that state no point estimate of the value, discussed
in \S\ref{sec:prior}.

That leaves 129 final worst-of contingent-coupon notes with a disclosed
estimated value, of which 71 can be priced from public option quotes. Of the 58
that cannot, 16 are written on single names or funds and 42 carry at least one
leg whose options are not listed. That second group is more concentrated than
one would guess. In 33 of the 42 the only such leg is the Nasdaq-100 Technology
Sector index, a US index with no options market of its own; the remaining nine
reference the EURO STOXX 50, the Nikkei 225 or a futures-based Russell variant.
The ratio is worth carrying forward. An outside party attempting this
verification can reach 55 per cent of the qualifying market, and what it cannot
reach is not a random sample of the rest but is dominated by notes on a
single index that has no options market of its own.

\subsection{The six notes}

% Generated by scripts/make_paper_tables.py from the frozen results.
% Do not edit by hand; the next build overwrites it.
\begin{table}[tbp]
\centering
\caption{The six notes: underlyings, terms and the figures the issuer discloses}
\label{tab:terms}
\small
\begin{tabular}{l l l r r r r r r}
\toprule
Note & Issuer & Underlyings & Tenor & Coupon & Barrier & Trigger & Est.\ value & Commission \\
\midrule
N1 & JPMorgan & NDX, RUT, SPX & 5.00 & 8.30 & 70.00 & 65.00 & 938.90 & 41.25 \\
N2\textsuperscript{a} & Morgan Stanley & DJX, NDX, RUT & 5.00 & 8.30 & 70.00 & 65.00 & 943.10 & 36.25 \\
N3 & Morgan Stanley & NDX, RUT, SPX & 3.01 & 9.75 & 70.00 & 60.00 & 983.40 & 2.50 \\
N4 & Goldman & NDX, RUT, SPX & 1.92 & 10.15 & 65.00 & 60.00 & 999.00 & 5.25 \\
N5 & BofA & NDX, RUT, SPX & 1.75 & 6.25 & 75.00 & 70.00 & 972.10 & 25.00 \\
N6 & JPMorgan & RUT, SPX, SMH & 1.92 & 12.40 & 65.00 & 60.00 & 967.30 & 22.25 \\
\bottomrule
\end{tabular}
\notetext{\textit{Notes:} Tenor in years from pricing date to stated maturity. Coupon is the contingent rate per annum in per cent, paid only when every underlying is at or above its coupon barrier on an observation date. \textit{Barrier} is that coupon threshold and \textit{Trigger} the principal-repayment threshold, both as a percentage of the initial level. Estimated value and selling commission are per \$1,000 of principal, as disclosed in the pricing supplement. \textsuperscript{a}\,Held out of the ranges quoted in the text: no fitted slice past one year carries a quoted strike at this note's trigger on one leg, so its downside is extrapolated rather than observed. Source: the filings listed in Appendix~\ref{app:terms}.}
\end{table}


Six of the 71 are priced here, listed in Table~\ref{tab:terms}, chosen to span
the structural variation rather than to be representative in any statistical
sense. They come from four issuers. Tenors run from 1.75 to 5.00 years, contingent
coupons from 6.25 to 12.40 per cent per annum, coupon barriers from 65 to 75
per cent of the initial level and downside triggers from 60 to 70 per cent. The
disclosed estimated values run from \$938.90 to \$999.00 per \$1{,}000 and the
selling commissions from \$2.50 to \$41.25, the low figure belonging to a note
sold only into fee-based advisory accounts where the adviser is compensated
separately.

The call mechanics vary more than the name of the product suggests, and the
variation is the reason six notes are enough to be interesting. Three are
redeemable at the issuer's option on scheduled dates. Two redeem under the
disclosed rational-exercise rule quoted in \S\ref{sec:intro}. One, and only
one, is auto-callable in the ordinary sense: redemption is mechanical when every
underlying closes at or above 93.25 per cent of its starting level on one of
six observation dates. That threshold is itself worth noting, since discussions
of these products often assume the automatic call sits at the initial level.

\subsection{Option data and the fitted surfaces}

Implied volatility surfaces are fitted to delayed end-of-day option chains
retrieved from Cboe's public endpoint and preserved as of 9 October 2026, with a
SHA-256 digest recorded for every file. Five surfaces cover the eighteen leg positions the six
notes carry between them, which resolve to five distinct underlyings. Expiries inside a fortnight are dropped, as are quotes whose
relative bid-ask exceeds 0.35, and log-moneyness is clipped to $[-0.9, 0.45]$;
consecutive expiries closer than 25 days apart are thinned so that the calendar
condition is tested across separated maturities rather than across
adjacent weeklies.

Each slice's forward and discount factor come from a put-call parity regression
across near-the-money strikes. Three of the five surfaces are fitted to cash
index options, which are European, and for these both quantities are fitted.
The other two are fitted to exchange-traded fund options, which are American,
and there the discount factor is pinned to the Treasury curve and only the
forward is fitted. Parity is an identity only under European exercise: left to
fit freely on the semiconductor ETF, the regression returned a discount factor
implying a 1.6 per cent zero rate against 4.3 per cent Treasuries, because the
early-exercise premium on in-the-money puts grows with strike and tilts the
slope. Pinning the discount factor removes that tilt while leaving the forward
to carry the market's dividend assumption. Near the money the early-exercise
premium itself is small enough to leave in the quotes, which is an
approximation rather than a correction and is revisited in \S\ref{sec:limits}.

\subsection{The one proxied leg}

Four of the five underlyings are priced off options on the instrument named in
the term sheet. The exception is the Dow Jones Industrial Average, whose listed
index options are too sparse to fit a surface across the maturities these notes
need, and which is therefore priced off options on the ETF that tracks it.

What that substitution costs can be measured rather than assumed, because two
other legs in the sample have both an index and an ETF chain available at shared
expiries. At-the-money implied volatility is higher on the fund than on the
index in every comparison: by 0.517 volatility points for the Dow pair at the
one expiry they share, and by 0.434 and 0.435 points on average for the
Nasdaq-100 and Russell 2000 pairs across 16 shared expiries each, reaching
beyond one and two years respectively. The sign is consistent and the magnitude
is similar across all three, which is what makes the Dow figure usable as an
estimate rather than a single observation.

Correcting that one leg, and nothing else, shifts the affected note's value from
\$893.93 to \$896.54, a rise of \$2.61 with a standard error of \$0.73, and
narrows its difference from the disclosed value from \$49.17 to \$46.56. The
proxy therefore accounts for about five per cent of that note's gap and moves it
in the direction that would close it. It is the same note whose downside trigger
falls outside the quoted strike range on the Dow leg, which is why
\S\ref{sec:results} reports it but holds it out of the stated ranges;
Appendix~\ref{app:coverage} gives the coverage counts leg by leg.

\subsection{Rates, spreads and the volatility indices}

Discounting uses the US Treasury par yield curve for the valuation date, as
published by the Treasury. The benchmark against which an implied funding spread
is judged is the ICE BofA A-rated US corporate index option-adjusted spread,
69 basis points on the pricing date, taken from the Federal Reserve Bank of St
Louis. The Cboe volatility indices for the S\&P 500, Nasdaq-100 and Russell
2000 supply the surface move between the pricing date and the snapshot date used
in \S\ref{sec:inputs}, each index matched to the leg it measures. Daily closes
for the semiconductor ETF come from a commercial endpoint rather than an
official series, no free official source for a single fund's history being
available; that series affects one note of the six.

Every series above is free and unauthenticated, and the retrieval code is in the
repository. The option chains the reported results depend on are preserved with
their digests, since a delayed-quote endpoint serves only current quotes and the
9 October snapshot cannot be retrieved again.
